\section{自主学习：已经发生的渐变}

\subsection{自主学习的多重含义}

``自主学习''在硅谷已成为热门话题，但每个人的定义不同。姚顺雨将其分为多种类型：

\begin{itemize}
\item \textbf{聊天个性化}：ChatGPT利用用户数据不断学习对话风格
\item \textbf{代码环境适应}：越来越熟悉每个公司独特的环境和文档
\item \textbf{科学探索}：像博士生一样从零学习有机化学成为领域专家
\end{itemize}

\begin{importantbox}{关键洞察：瓶颈在数据/任务，不在方法论}
自主学习的真正瓶颈不是方法论（how），而是数据和任务（what and where）——在什么场景下、基于什么前提、用什么样的奖励函数去做。
\end{importantbox}

\subsection{自主学习已经在发生}

\begin{knowledgebox}{Claude Code的自我改进}
Claude Code已经写了自身项目95\%的代码——从某种意义上说，它在帮助自己变得更好。这不就是一种自主学习吗？

类似地，Cursor的AutoComplete Model每几小时就用最新用户数据训练，Composer Model也在使用真实环境数据。这些都是自主学习的早期信号。
\end{knowledgebox}

姚顺雨在2022--2023年做SWE-Agent时就提出：如果有一天SWE-Agent能自己改进SWE-Agent的代码仓库，那就是一种自我学习。Claude Code已经在大规模做这件事了。

\begin{warningbox}{渐变而非突变}
姚顺雨认为自主学习更像渐变而非突变——目前的例子还局限在特定场景下，但趋势是明确的。最大的瓶颈是\textbf{想象力}：如果2026--2027年有新范式出现，那个Blog Post应该长什么样？它应该展示什么样的任务和效果？
\end{warningbox}


